\section{Future directions}
\label{sec:future}

The library is intentionally a foundation:\ it is meant to be
extended.  We list the extensions we view as natural, in
descending order of how much existing infrastructure they would
re-use.

\subsection{Algorithmic and constructive Nash}
\label{sec:future:algo}

\Cref{thm:nash-mixed} is non-constructive (Brouwer); a natural
extension is a constructive proof for two-player games via the
\emph{Lemke–Howson} pivoting algorithm~\citep{LemkeHowson1964} or
via \emph{Sperner's lemma}~\citep{Sperner1928} (a combinatorial
fixed-point statement on triangulated simplices that yields
Brouwer constructively for low dimensions), packaged as an
algorithmic witness function on top of the existing strategic-form
definitions.  The \emph{PPAD}-hardness
story~\citep{DaskalakisGoldbergPapadimitriou2009} — where PPAD is
the complexity class of problems whose existence is guaranteed by
a parity argument on a directed graph, the canonical example being
end-of-line — is at a different level (it concerns problem
reductions, not proof-of-existence) and is largely orthogonal:\ a
complexity formalization would need an independent computation
model.

\subsection{Sequence form and efficient EFG solving}
\label{sec:future:sequence-form}

A \emph{sequence} in an extensive-form game with perfect recall is
the list of own-actions a player has taken along a history.  The
\emph{sequence-form representation}
of~\citet{KollerMegiddoVonStengel1994} (see also the survey
of~\citet{vonStengel2002}) replaces the exponential-sized pure
strategy space by realization probabilities indexed by sequences,
turning equilibrium computation into a linear or quadratic
program of size polynomial in the tree.  Adding the sequence-form bridge as a fourth EFG
strategy carrier (alongside pure, behavioral, mixed) would let
the library state and prove correctness of the standard solving
algorithms (LP duality$\Rightarrow$two-player zero-sum value;
sequence-form quadratic program for general-sum games).  The
Kuhn-theorem infrastructure of \Cref{sec:theorems:kuhn} already
supplies the recall layer.

\subsection{Counterfactual regret minimization}
\label{sec:future:cfr}

\emph{Counterfactual regret minimization} (CFR,
\citealp{Zinkevich2007CFR}) is the dominant family of algorithms
for finding approximate equilibria in large imperfect-information
extensive-form games; it works by tracking, at each information
set, the cumulative regret for not having played each action,
counterfactually weighted by the probability of reaching that
information set under the opponents' play.  The regret-equilibrium
duality and the no-regret learning machinery a CFR proof would
consume are already formalized (\Cref{sec:solution-concepts:regret},
\Cref{sec:theorems:learning}) — including explicit routes to
$\varepsilon$-CCE via no-regret self-play without going through
Brouwer at each finite horizon.  What remains specific to CFR is the counterfactual,
information-set-local regret decomposition for extensive-form games
and its convergence proof on top of that base.

\subsection{Imperfect-recall extensions}
\label{sec:future:absent-mind}

The Kuhn theorem in the library handles a recall hypothesis
weaker than perfect recall (per-step ActionPosteriorLocal).  The
basic \emph{absent-minded-driver}
phenomenon~\citep{PiccioneRubinstein1997} is formalized in
\TT{Languages/MultiRound/AbsentMindedDriver.lean}.  The protocol
has one player and two indistinguishable decision rounds; it
violates \TT{ViewDeterminesRound}, the MultiRound
no-absent-mindedness condition.  The payoff calculation is also
mechanized:\ with payoffs $0$ for the first exit, $4$ for the
second exit, and $1$ for missing both exits, the ex-ante optimal
continuation probability is $2/3$, while decision-time
re-optimization at the single information set under those arrival
odds chooses $1/3$ and strictly improves the decision-time value.

Natural extensions are the broader imperfect-recall pathologies
around this example:\ general conditions for time consistency,
relations between information-set revisitation and behavioral
optimality, and representation-neutral EFG/FOSG versions of the
same construction.  These should live on the observation-model
layer already used by Kuhn's theorem, with the MultiRound example
serving as the finite sanity check.

\subsection{Infinite and continuous games}
\label{sec:future:infinite}

The largest single extension would be to relax the discrete-PMF
discipline and use mathlib's measure theory.  This would let the
library state and prove:\ Glicksberg's theorem on continuous
payoff functions over compact strategy
spaces~\citep{Glicksberg1952}, the Milgrom–Weber existence theorem
for Bayesian games with uncountable type
spaces~\citep{MilgromWeber1985}, infinite-horizon discounted
stochastic games~\citep{Shapley1953stochastic, MertensNeyman1981},
and continuous-action or public-randomization variants of the folk
theorem.  The current library proves the discrete observable
mixed-action approximate discounted folk theorem of
\Cref{thm:folk-discounted}; a measure-theoretic extension would
give it continuous companions.

This is a substantial extension and would change the library's
character.  The fixed-point ingredient is partly isolated:
\TT{Math/SchauderFixedPoint.lean} proves Schauder's theorem for
compact convex subsets of real normed spaces.  What remains for
continuous games is not just a fixed-point theorem but the
measure-theoretic and topological game layer around it.  We see
that less as a single project and more as a parallel library with
shared infrastructure (the
\TT{KernelGame}/\TT{GameForm} skeleton generalizes to a kernel
$\K : \prod_i \Strat_i \to \mathcal{P}(\Outcome)$ for a measure
type $\mathcal{P}$, with the discrete \TT{PMF} as one
instantiation).

\subsection{Cooperative game theory}
\label{sec:future:cooperative}

The cooperative branch lives at \TT{Cooperative/} (a parallel
paradigm to the non-cooperative core occupying the rest of the
library; see \Cref{sec:theorems:shapley-uniqueness} for the
framing).  Already formalized:\ TU coalitional games with the
Shapley value~\citep{Shapley1953} and the four classical axioms
used by the Lean uniqueness theorem,
the unanimity-basis decomposition (Möbius inversion on the subset
lattice), the full Shapley uniqueness theorem
(\Cref{thm:shapley-unique}), the Banzhaf~\citep{Banzhaf1965} and
Shapley–Shubik~\citep{ShapleyShubik1954} power indices, simple
(voting) games, convex (supermodular) games with the
$\mathrm{IsConvex} \Leftrightarrow$ monotone-marginals characterization
\citep{Shapley1971} and closure under sum and nonnegative scalar,
the core (efficient $\land$ coalition-rational) with the
convex-set closure property and individual rationality, core
closure under finite sums and nonnegative scalar multiplication,
the Shapley-in-core theorem for convex games (hence nonempty core),
the positive-cone specialization through nonnegative
unanimity-basis coefficients, additive games (the identity
allocation), and the 3-player majority game as a worked example.
The \emph{cost of stability}~\citep{Bachrach2009} — the least external
subsidy that makes the grand coalition stable — is formalized
(\TT{IsStabilizable}, \TT{costOfStability}) with its monotonicity, the
vanishing of the cost on games with a nonempty core, and zero-cost
classes (nonnegative unanimity coefficients, additive games).
Balanced collections and the easy direction of the
\emph{Bondareva–Shapley} theorem~\citep{Bondareva1963, Shapley1967}
— a nonempty core implies balancedness — are formalized
(\TT{IsBalancedCollection}, \TT{CoalGame.IsBalanced},
\TT{IsCore.isBalanced}) by the double-counting swap over a core
allocation.
Nash bargaining
\citep{Nash1950bargaining} is formalized through the IR / weak-Pareto
characterization (the Nash maximizer is only weakly Pareto on the
disagreement boundary in this general non-convex setting), with
symmetric and affine-invariant variants;
the \emph{egalitarian} bargaining solution~\citep{Kalai1977} is also
formalized as \TT{IsEgalitarian} (the maximal feasible equal-gain
point), together with the fact that it gives at least the common gain
of the symmetric Nash solution, and the
\emph{Kalai–Smorodinsky} solution~\citep{KalaiSmorodinsky1975} through
its proportionality characterization
(\TT{IsKalaiSmorodinskyRelativeTo}, carrying IR, strong Pareto,
symmetry, and affine invariance for any reference point), with the
solution proper (\TT{IsKalaiSmorodinsky}) fixing the ideal point and
bounding each player's gain by it.
Stable matching~\citep{GaleShapley1962} is formalized through the
\TT{IsMatching}/\TT{IsBlockingPair}/\TT{IsStable} predicates, with the
existence of a stable matching proved via the Gale–Shapley
deferred-acceptance algorithm (\Cref{thm:gale-shapley}).

Natural extensions:\ the converse (hard) direction of
Bondareva–Shapley \citep{Bondareva1963, Shapley1967} — balancedness
implies a nonempty core — which closes the characterization of
core-nonempty games and is the flagship LP-duality (Farkas) result;
the nucleolus \citep{Schmeidler1969}; the Aumann–Shapley value for non-atomic
games~\citep{AumannShapley1974}; NTU coalitional games; and further
bargaining solutions.

\subsection{Information design and signaling}
\label{sec:future:info-design}

In \emph{Bayesian persuasion}~\citep{KamenicaGentzkow2011} a
sender designs an information structure (a signal-generation
mechanism that maps the state of the world to a public message)
to influence a receiver's posterior beliefs and hence their
action; the design problem is to choose the most persuasive
signal mechanism subject to the receiver's Bayesian rationality.
\emph{Signaling games}~\citep{Spence1973} are the strategic
analogue with an informed sender choosing a costly signal whose
cost may itself depend on type.  \emph{Information
design}~\citep{Bergemann2019} generalizes the persuasion setup to
multiple receivers and richer outcome rules.

The finite substrate is formalized in
\TT{Mechanism/InformationDesign.lean}:\ finite signal structures
$\Omega \to \PMF{M}$, induced joint state-message laws, the
message marginal formula, the Bayes-plausibility invariant,
uninformative and full-information signals, finite sender/receiver
persuasion problems, persuasive decision rules, and finite-sum
expressions for sender expected utility.  This is the right layer
for discrete persuasion examples and for connecting disclosure
problems back to the Bayesian-game infrastructure of
\Cref{sec:mech:bayesian}.  A complementary belief-based layer
(\Cref{sec:mech:info-design}) provides the Bayes-plausibility
\emph{splitting} characterization via the canonical posterior coupling
and its multi-receiver necessary direction~\citep{Arieli2021}.

Natural extensions are substantive theorems over this substrate:
finite concavification of the optimal sender value, value-of-information
results, optimal signal existence on a bounded finite message alphabet,
explicit finite signaling-game refinements, positive characterizations
or counterexample libraries for multi-receiver feasibility, and bridges
from persuasive decision rules to direct mechanisms.  Continuous type spaces and compact belief simplices
would require the measure-theoretic extension discussed in
\Cref{sec:future:infinite}; the finite layer is useful independently
and should remain separate from that analytic generalization.

\subsection{Social choice impossibility theorems}
\label{sec:future:social-choice}

\emph{Arrow's impossibility theorem}~\citep{Arrow1951} — no
social-welfare function aggregating individual ordinal preferences
over $\ge 3$ alternatives can simultaneously satisfy unanimity,
non-dictatorship, and independence of irrelevant alternatives — is
formalized (\Cref{sec:mech:social}, \TT{Mechanism/Arrow.lean}),
via Geanakoplos's pivotal-voter argument on a strict-linear-order
rationality layer.  Its strategic analogue
\emph{Gibbard–Satterthwaite} \citep{Gibbard1973, Satterthwaite1975}
— any non-dictatorial social choice function with full range over
$\ge 3$ alternatives is manipulable — is formalized as well
(\TT{Mechanism/GibbardSatterthwaite.lean}), by the reduction to
Arrow through an induced social-welfare function.
\emph{Sen's liberal paradox}~\citep{Sen1970}, showing that even a
minimal individual-rights condition is inconsistent with Pareto
efficiency plus universal domain, is formalized as well
(\Cref{sec:mech:social}, \TT{Mechanism/SenParetianLiberal.lean}), on
the same strict-ranking social-choice layer.  With the ordinal
impossibility results in place, the natural extensions are the
randomized and quantitative refinements — probabilistic social choice
and manipulation-robustness bounds — that would build on the library's
discrete-probability layer.

\subsection{Multi-agent learning}
\label{sec:future:learning}

With the no-regret core in place (\Cref{sec:theorems:learning}), the
open directions are the asymptotic and internal-regret refinements.
The multiplicative-weights and self-play bounds are stated for a
\emph{fixed} learning rate, where the residual error does not
vanish; the horizon-dependent tuning $\eta \approx \sqrt{\log|A|/T}$
that makes external regret sublinear — and so turns the per-horizon
$\varepsilon$-bound into a genuine \emph{convergence} statement — is
the natural next step \citep{Cesa-BianchiLugosi2006,
BlumMansour2007}.  Beyond it:\ internal-regret minimization driving
play to \emph{correlated} (not merely coarse correlated)
equilibrium; Robinson's theorem that fictitious play converges in
two-player zero-sum games; and the reinforcement-learning algorithms
(self-play with function approximation, MCTS) that the FOSG and
observation-model layers were built to host.

\subsection{Stable matching beyond existence}
\label{sec:future:matching}

The Gale–Shapley deferred-acceptance algorithm and the existence of a
stable matching from an arbitrary strict preference
profile~\citep{GaleShapley1962} are formalized
(\Cref{thm:gale-shapley}, \TT{Cooperative/GaleShapley.lean}), together
with proposer-optimality and receiver-pessimality
(\TT{daMatching\_man\_optimal}, \TT{daMatching\_woman\_pessimal}).  The
natural extensions are the remaining structural results:\ the lattice
structure of the set of stable matchings, strategy-proofness for the
proposing side, and the many-to-one (hospitals/residents)
generalization with the rural-hospitals theorem.  A computable
extraction of the $O(n^2)$ algorithm paired with the existing
correctness proof is a further refinement; Bertot's Coq
formalization \citep{Bertot2016} is a useful reference point.

\subsection{Type-theoretic refactoring opportunities}
\label{sec:future:refactor}

We close with two design directions internal to the library
itself.  First, an open-games-style
reformulation~\citep{GhaniHedgesWinschelZahn2018} would expose the
\TT{KernelGame} category in a shape suitable for compositional
reasoning at a higher abstraction level; this would not change
the theorems but might simplify their proofs.  Second, the
library's \TT{noncomputable} discipline could be tightened by
isolating the constructive fragment (definitions and theorems
that do not invoke classical reasoning) as a separately-named
namespace; this would make the boundary between the constructive
finite-game fragment and the non-constructive existence-of-Nash
fragment visible in user code.
